\uuid{ACSo}
\titre{Variables continues}
\auteur{worms}
\organisation{exo7}
\licence{CC BY-NC-SA 4.0}
\datecreate{2017-08-28}
\niveau{L2}
\module{Probabilité et statistique}
\chapitre{Probabilité continue}
\sousChapitre{Densité de probabilité}
\difficulte{}
\theme{Variables continues}
\exo7qcmid{worms.worms.044}

\contenu{
\begin{question}
\qclassification{261.01}
Soit $X$ une variable aléatoire de loi ${\cal N}(30,25)$. Lesquelles (ou laquelle) des affirmations suivantes sont vraies ?
\begin{answers}
\bad{$X$ a environ 5\% de chances d'être entre $-20$ et $80$, et $\mathbb{P}(X>30)=1/2$.}
\good{Le premier quartile de $X$ est environ égal à $26.7$ et $\mathbb{P}(X<30)=1/2$.}
\bad{La loi de $Y=2X$ est ${\cal N}(60,50)$ et $\mathbb{P}(20\leq X\leq 40)\simeq 95\%$.}
\good{$\mathbb{P}(|X-30|>20)\simeq 0$ et la loi de $Y=2X+10$ est ${\cal N}(70,100)$.}
\end{answers}
\begin{explanations}
Dans cette question, il faut surtout prendre garde à ne pas confondre la variance et l'écart-type.
\end{explanations}
\end{question}
}
